\subsection{应用：I. 有理数及一元有理函数的典范分解}

定理2.——设A为主理想整环，K为其分式域，P为A中不可约元素的代表元系。给定一个元素 \(x \in K\) ，存在 P 的一个有限子集 H、元素 \(a_{0} \in A\) 与 \(a_{p} \in A\) （在 A 中不被 p 整除， \(p \in H\) ）、以及整数 \(s(p) > 0\) ( \(p \in H\) ）使得

\[
x = a _ {0} + \sum_ {p \in \mathrm{H}} a _ {p} p ^ {- s (p)},\tag{1}
\]

其中H和 \(s(p)\) 由这些条件唯一确定。

此外，如果 \(R_{p}\) 表示 A 的一个恰好包含 A 中每个模 p 剩余类中一个元素的子集（ \(p \in P\) ），则每个 \(x \in K\) 可以唯一地表示为

\[
x = a + \sum_ {p \in P} \left(\sum_ {h = 1} ^ {\infty} r _ {p h} p ^ {- h}\right)\tag{2}
\]

，其中 \(a \in \mathbf{A}\) ，且 \(r_{ph} \in \mathbf{R}_p\) 对所有 \(h\) 与 \(p\) 成立，并且除有限多个外 \(r_{ph}\) 均为零。

将 K 视为 A-模；则 A 是由 1 生成的子模。商模 K/A 是 K 关于等价关系 \(x' - x \in \mathbf{A}\) 的商，在 VI 第 6 页的记号中也写作 \(x \equiv x' \pmod{1}\) ；设 \(f\) 表示从 K 到 M = K/A 上的自然同态。

商模 M 是一个挠模，因为 M 的每个元素都具有形式 \(f(a/b)\) ( \(a \in A\) , \(b \in A\) , \(b \neq 0\) ），从而 \(bf(a/b) = f(a) = 0\) 。因此，VII 第 8 页的定理 1 适用。设 \(M_{p}\) ( \(p \in P\) ) 为 M 中被 p 的幂零化的元素构成的子模；则 \(f^{-1}(M_{p})\) 是子环 \(A_{p}\) （由 K 中形如 \(ap^{-n}\) 的元素组成，其中 \(a \in A\) ， \(n \geqslant 0\) 为整数）。模 M 是这些 \(M_{p}\) 的直和，因此每个 \(x \in K\) 在模 1 下同余于这些 \(A_{p}\) 之和中的一个元素；换言之，公式 (1) 成立，其中 \(s(p)\) 为 \textgreater0 的整数，且诸 \(a_{p}\) 是 A 中有限多个元素，使得 \(a_{p}\) 不是 p 的倍数。

我们现在证明这些关于 \(s(p)\) 与 \(a_{p}\) 的条件完全确定了 H 和 \(s(p)\) 。事实上，H 就是 \(p \in \mathbf{P}\) 中使得 \(f(x)\) 在 \(\mathbf{M}_{p}\) 中的分量 \(\neq 0\) 的那些 p 的集合。另一方面，如果 \(s\) 与 \(s'\) 是两个 \(>0\) 的整数，满足 \(s \geqslant s'\) ，并且 \(a\) 与 \(a'\) 是 \(A\) 的不被 \(p\) 整除的元素，满足 \(ap^{-s} \equiv a'p^{-s'} \pmod{1}\) ，那么我们推断出 \(a \equiv a'p^{s - s'} \pmod{p^s}\) ；如果 \(s > s'\) ，则 \(a \equiv 0 \pmod{p}\) ，这与假设相矛盾。这个论证还表明每个 \(a_p\) 在模 \(p^{s(p)}\) 下是唯一确定的.

为了完成证明，我们首先注意到，在每个模 \(p^s\) 的剩余类（在 \(\mathbf{A}\) 中）里，存在唯一一个形如 \(\sum_{h=0}^{s-1} r_h p^h\) 的元素，其中 \(r_h \in \mathbb{R}_p\) ， \(0 \leqslant h \leqslant s-1\) 。事实上，我们对 \(s\) 用归纳法（该性质由 \(\mathbb{R}_p\) 的定义在 \(s=1\) 时得出）：设 \(x \in \mathbf{A}\) ；根据假设，存在唯一一个形如 \(\sum_{h=0}^{s-2} r_h p^h\) ( \(r_h \in \mathbb{R}_p\) ) 的元素，属于 \(x \mod p^{s-1}\) 的剩余类；则差 \(x - \sum_{h=0}^{s-2} r_h p^h\) 形如 \(ap^{s-1}\) ，即 \(p^{s-1}\) 的一个倍数；现在存在唯一的元素 \(r_{s-1}\) （ \(\mathbb{R}_p\) 的）使得 \(a \equiv r_{s-1} (\bmod p)\) ；由此 \(x \equiv \sum_{h=0}^{s-1} r_h p^h (\bmod p^s)\) 。为了得到公式 (2)，现在只需将此事实应用于公式 (1) 中的每个 \(a_p\) 。鉴于上述，唯一性是显然的。

以下是定理2最重要的应用：

I. 环 A 是有理整数环 Z，且 K = Q。设 P 为 \textgreater{} 0 的素数集合，且对于每个 \(p \in P\) ，令 \(R_{p}\) 为 Z 中的区间 \((0, p - 1)\) 。因此我们有典范分解

\[
x = a + \sum_ {p \in P} \left(\sum_ {h = 1} ^ {\infty} e _ {p h} p ^ {- h}\right)
\]

，其中 \(a \in \mathbf{Z}\) , \(e_{ph} \in \mathbf{Z}\) ，且 \(0 \leqslant e_{ph} \leqslant p - 1\) .

\begin{enumerate}
\def\labelenumi{\Roman{enumi}.}
\setcounter{enumi}{1}
\tightlist
\item
  环 A 是交换域 E 上一个未定元的多项式环 E{[}X{]}，且 \(\mathbf{K} = \mathbf{E}(\mathbf{X})\) 。设 P 为 E{[}X{]} 中首一不可约多项式的集合（VII，第 5 页）。对于 \(p \in P\) ，根据多项式的欧几里得除法（IV，第 10 页），我们可以取 \(R_{p}\) 为次数严格小于 p 的次数的多项式集合。因此我们得到有理函数 \(r(\mathbf{X}) \in \mathbf{E}(\mathbf{X})\) 的（称为典范分解的）分解：
\end{enumerate}

\[
r (\mathrm{X}) = a (\mathrm{X}) + \sum_ {p \in \mathrm{P}} \left(\sum_ {h = 1} ^ {\infty} v _ {p h} (\mathrm{X}) \cdot p (\mathrm{X}) ^ {- h}\right)
\]

，其中 \(a(\mathbf{X})\) 是一个多项式，且 \(v_{ph}(\mathbf{X})\) 是次数严格小于 \(p(\mathbf{X})\) 的次数的多项式（对所有 \(p\) 与 \(h\) ）。特别地，如果域 \(\mathbf{E}\) 是代数闭的，则诸 \(p(\mathbf{X})\) 具有形式 \(\mathbf{X} - \alpha\) （ \(\alpha \in \mathbf{E}\) ），从而诸 \(v_{ph}(\mathbf{X})\) 是常数。

因此我们可以说，域 E 上的向量空间 \(E(X)\) 具有一个基，由单项式 \(X^{n}\) ( \(n \geqslant 0\) 为整数）以及形如 \(\mathbf{X}^{m}/(p(\mathbf{X}))^{h}\) 的有理函数组成，其中 \(p \in P\) ，h 和 m 是满足 \(h \geqslant 1\) 与 \(0 \leqslant m < \deg(p)\) 的整数.

